\section{Weighted curvature and stable surfaces}
\label{sec:weighted-surfaces}

The spectral hypothesis produces a positive density for which closed
stable weighted minimal surfaces have uniformly bounded diameter.  We
prove both the closed-surface estimate and a local version that does not require the surface to be
complete.  The latter will be used to extract compact surfaces in
Section~\ref{sec:cutting}.

We use the sign convention $\Delta=\Div\nabla$ and take round spheres to
have positive scalar curvature.  A closed surface is compact without
boundary.  All integrals use Riemannian densities, so none of the formulas
requires an orientation.  For a Riemannian metric
$h$ and a smooth real function $T$ on a manifold of any dimension, set
\begin{equation}\label{eq:weighted-scalar}
 D(h,T)=\Scal_h+2\Delta_hT-|dT|_h^2.
\end{equation}
We work with the normalization $\lambda=1$ until the final rescaling in
Section~\ref{sec:graph-gluing}.

\subsection{From the quadratic form to a positive density}

We use the standard ground-state construction from a nonnegative
quadratic form; see Fischer-Colbrie--Schoen~\cite[Theorem~1]{FischerColbrieSchoen1980}.
We recall the exhaustion argument in the form needed here.

\begin{lemma}[Positive weight]\label{lem:positive-weight}
Let $(M,g)$ be a connected smooth Riemannian manifold without boundary.
Suppose that
\[
 \int_M\bigl(4|d\phi|_g^2+\Scal_g\phi^2\bigr)\,d\vol_g
 \ge \int_M\phi^2\,d\vol_g
 \qquad(\phi\in C_c^\infty(M;\R)).
\]
There is a smooth function $t:M\to\R$ such that $D(g,t)\ge1$.
\end{lemma}

\begin{proof}
It suffices to find a positive smooth function $v$ satisfying
\[
 (-4\Delta_g+\Scal_g)v\ge v,
\]
since $t=-2\log v$ then gives
\begin{equation}\label{eq:log-weight}
 D(g,t)=\Scal_g-4\frac{\Delta_gv}{v}\ge1.
\end{equation}
If $M$ is compact, take a positive first eigenfunction of
$-4\Delta_g+\Scal_g$.  Its eigenvalue is at least one by the hypothesis.

If $M$ is noncompact, fix $o\in M$ and choose an increasing exhaustion
$\Omega_j$ by connected relatively compact smooth domains containing
$o$.  Let $\lambda_j$ and $v_j$ be the first Dirichlet eigenvalue and a
positive first eigenfunction on $\Omega_j$, normalized by $v_j(o)=1$.
The quadratic-form hypothesis, extended to $H_0^1(\Omega_j)$ by density,
gives $\lambda_j\ge1$.  Domain monotonicity makes $\lambda_j$ decreasing,
and a fixed nonzero test function in $\Omega_1$ bounds it from above.
Thus $\lambda_j\to\lambda_\infty\ge1$.  On each compact subset of $M$,
the local Harnack inequality, applied along finitely many chains of
coordinate balls joining it to $o$, bounds $v_j$ above and away from zero
for all sufficiently large $j$.  Interior elliptic estimates
\cite[Theorems~6.2 and~8.20]{GilbargTrudinger2001} and a
diagonal subsequence now give smooth convergence on compact subsets to
a positive function $v$ with
\[
 (-4\Delta_g+\Scal_g)v=\lambda_\infty v.
\]
Equation~\eqref{eq:log-weight} proves the assertion.
\end{proof}

Fix this function $t$ on $M$.  Its density $e^{-t}$ determines the weighted
area used to cut $M$.  We first carry out the surface calculation for
arbitrary ambient data $(h,T)$, since a modified metric and weight will
also occur later.

\subsection{The stability inequality}

Let $\Sigma$ be a smooth two-sided surface, possibly with boundary, in a
Riemannian three-manifold $(X,h)$, and choose its unit normal $\nu$.
All surface variations and test functions will be compactly supported
in $\operatorname{int}\Sigma$.  Write $k=h|_\Sigma$ and use
the conventions
\[
 B(V,W)=h(\nabla_V\nu,W),\qquad H=\operatorname{tr}_k B.
\]
For a smooth ambient function $T$, weighted area means
$\int_\Sigma e^{-T}\,d\Area_k$.  If $\Sigma$ is the boundary of a set
$E$, with outward normal $\nu$, the functional
\begin{equation}\label{eq:weighted-pressure-functional}
 \int_{\partial E}e^{-T}\,d\Area_h
 -\int_E\mu e^{-T}\,d\vol_h
\end{equation}
has Euler--Lagrange equation
\begin{equation}\label{eq:weighted-mean-curvature}
 H-\nu T=\mu,
\end{equation}
where $\mu$ is a smooth function on $X$.  For a surface not globally
bounding a set, the same equation and second variation are defined
locally using its chosen normal.  Stability always means unconstrained
stability: all normal variations compactly supported in the surface
interior are allowed, with no volume constraint.

The weighted surface Laplacian and the potential in the second variation
are
\begin{equation}\label{eq:weighted-jacobi}
 \begin{split}
 \Delta_{\Sigma,T}u
   &=\Delta_ku-\langle d(T|_\Sigma),du\rangle_k,\\
 Q&=\Ric_h(\nu,\nu)+\Hess_hT(\nu,\nu)+|B|_k^2+\nu\mu.
 \end{split}
\end{equation}
In particular, stability is the inequality
\begin{equation}\label{eq:stability-form}
 \int_\Sigma\bigl(|du|_k^2-Qu^2\bigr)e^{-T}\,d\Area_k\ge0
 \qquad(u\in C_c^\infty(\operatorname{int}\Sigma)).
\end{equation}
To check the signs, a normal variation of speed $u$ satisfies
\[
 H'=-\Delta_ku-(\Ric_h(\nu,\nu)+|B|_k^2)u,
 \qquad \nu'=-\nabla_ku.
\]
Differentiating $H-\nu T-\mu$ therefore gives
$-\Delta_{\Sigma,T}u-Qu$.  The first variation of
\eqref{eq:weighted-pressure-functional} has factor $H-\nu T-\mu$,
so its second variation at a stationary surface is exactly
\eqref{eq:stability-form}.

\begin{lemma}[Curvature transferred to a stable surface]
\label{lem:weighted-stability}
Suppose that a smooth two-sided surface $\Sigma$ satisfies
\eqref{eq:weighted-mean-curvature} and is stable in the sense of
\eqref{eq:stability-form}.  Let $\Omega$ be either the whole surface,
if it is closed and connected, or a connected smooth surface domain
with compact closure in $\operatorname{int}\Sigma$.  There is a smooth
positive function $w$ in $\Omega$ such that, with $p=\log w$ and
$U=T|_\Sigma-p$,
\begin{equation}\label{eq:surface-curvature-transfer}
 D(k,U)\ge
 D(h,T)+\mu^2+|B|_k^2+2\nu\mu+|dp|_k^2
 \qquad\text{in }\Omega.
\end{equation}
In the domain case, no completeness hypothesis is imposed on $\Sigma$.
\end{lemma}

\begin{proof}
Take a positive first eigenfunction of the self-adjoint operator
$-\Delta_{\Sigma,T}-Q$ in the measure $e^{-T}\,d\Area_k$, with
Dirichlet boundary condition in the domain case.  Stability makes its
first eigenvalue nonnegative.  Thus
\begin{equation}\label{eq:log-jacobi}
 -\Delta_{\Sigma,T}w\ge Qw,
 \qquad
 -\Delta_{\Sigma,T}p\ge Q+|dp|_k^2.
\end{equation}
The function $w$ is positive in $\Omega$; its possible vanishing on
$\partial\Omega$ will not enter the argument.

The Gauss equation and the decomposition of the ambient Laplacian read
\[
 \Scal_k=\Scal_h-2\Ric_h(\nu,\nu)+H^2-|B|_k^2,
 \qquad
 \Delta_hT=\Delta_k(T|_\Sigma)+H\nu T+\Hess_hT(\nu,\nu).
\]
Together with $|dT|_h^2=|d(T|_\Sigma)|_k^2+(\nu T)^2$, these give
\begin{equation}\label{eq:surface-gauss-weighted}
 \begin{split}
 D(k,T|_\Sigma)={}&D(h,T)
 -2\bigl(\Ric_h(\nu,\nu)+\Hess_hT(\nu,\nu)\bigr)\\
 &+(H-\nu T)^2-|B|_k^2.
 \end{split}
\end{equation}
Changing $T|_\Sigma$ to $T|_\Sigma-p$ gives the exact identity
\[
 D(k,U)=D(k,T|_\Sigma)-2\Delta_{\Sigma,T}p-|dp|_k^2.
\]
Now substitute \eqref{eq:log-jacobi},
\eqref{eq:weighted-jacobi}, and $H-\nu T=\mu$ into
\eqref{eq:surface-gauss-weighted}.  The ambient normal Ricci and Hessian
terms cancel, leaving \eqref{eq:surface-curvature-transfer}.
\end{proof}

In particular, the lower bound
\begin{equation}\label{eq:pressure-curvature-condition}
 D(h,T)+\mu^2-2|d\mu|_h\ge c>0
\end{equation}
along a stable surface gives $D(k,U)\ge c$ on the closed surface or
the chosen relatively compact domain in Lemma~\ref{lem:weighted-stability}.
We next convert this curvature inequality into a length estimate in the
original induced metric $k$, not merely in a conformal metric.

\subsection{The conformal length estimate}

The conformal geodesic argument is part of the stable-surface radius
method of Schoen--Yau~\cite{SchoenYau1983}; the differential inequality
used here also appears in Chodosh--Li~\cite[Lemma~16]{ChodoshLi2024}.
We give the index calculation explicitly to track length in the original
metric.

\begin{lemma}[Conformal index estimate]\label{lem:conformal-length}
Let $(S,k)$ be a smooth surface, let $U\in C^\infty(S)$, and suppose
that $D(k,U)\ge c>0$.  Let $\gamma$ be a finite geodesic segment for
$\widetilde k=e^{-2U}k$ whose index form is nonnegative for normal
variations compactly supported in its interior.  Then
\[
 \operatorname{length}_k(\gamma)\le\frac{2\pi}{\sqrt c}.
\]
The same conclusion holds when the curvature hypothesis is assumed only
on a neighborhood of the segment's interior.
\end{lemma}

\begin{proof}
The Gauss curvature of $\widetilde k$ is
\begin{equation}\label{eq:conformal-gauss}
 K_{\widetilde k}=e^{2U}(K_k+\Delta_kU)
 \ge\tfrac12 e^{2U}(c+|dU|_k^2),
\end{equation}
because $\Scal_k=2K_k$.  Parametrize $\gamma$ by $k$-arclength
$s\in[0,L]$.  Choose a parallel $\widetilde k$-unit normal along its
interior; such a choice exists along an interval even when $S$ is
nonorientable.  Since $d\widetilde s=e^{-U}ds$, the index inequality
and \eqref{eq:conformal-gauss} imply
\begin{equation}\label{eq:conformal-index}
 0\le\int_0^L e^U
 \left((\phi')^2-\tfrac12(c+|dU|_k^2)\phi^2\right)\,ds
 \qquad(\phi\in C_c^\infty(0,L)).
\end{equation}
Here and below $U$ and its differential are evaluated along $\gamma$.
Set $\phi=e^{-U/2}\psi$.  As $|U'|\le|dU|_k$, completing the square
gives
\begin{align*}
 e^U\left((\phi')^2-\tfrac12(c+|dU|_k^2)\phi^2\right)
 &=2(\psi')^2-\tfrac c2\psi^2
   -\left(\psi'+\tfrac12U'\psi\right)^2\\
 &\quad-\tfrac12\bigl(|dU|_k^2-(U')^2\bigr)\psi^2\\
 &\le2(\psi')^2-\tfrac c2\psi^2.
\end{align*}
Consequently
\[
 2\int_0^L(\psi')^2\,ds\ge\frac c2\int_0^L\psi^2\,ds
 \qquad(\psi\in C_c^\infty(0,L)).
\]
Approximating $\sin(\pi s/L)$ in $H_0^1(0,L)$ gives
$2\pi^2/L^2\ge c/2$, which is the asserted bound.  All test variations
are supported away from the endpoints, proving also the last assertion.
\end{proof}

\begin{proposition}[Closed stable surfaces]\label{prop:stable-surface}
Let $\Sigma$ be a smooth closed connected two-sided surface in a
Riemannian three-manifold $(X,h)$.  Let $T,\mu$ be smooth ambient
functions.  Suppose that $\Sigma$ satisfies $H-\nu T=\mu$, is stable
for weighted area with pressure $\mu$, and satisfies
\eqref{eq:pressure-curvature-condition} for a constant $c>0$.
Then, in the induced metric $k=h|_\Sigma$,
\[
 \diam_k\Sigma\le\frac{2\pi}{\sqrt c}.
\]
Moreover, $\Sigma$ is diffeomorphic to $S^2$ or $\R P^2$.
\end{proposition}

\begin{proof}
Lemma~\ref{lem:weighted-stability} gives a smooth function $U$ on
$\Sigma$ with $D(k,U)\ge c$.  The conformal metric $e^{-2U}k$ is
complete because $\Sigma$ is compact.  Join any two points by a
minimizing geodesic for that metric.  Its index form is nonnegative,
so Lemma~\ref{lem:conformal-length} bounds its $k$-length by
$2\pi/\sqrt c$.  The $k$-distance between the endpoints is no greater.

Integrating $D(k,U)\ge c$ and using Gauss--Bonnet gives
\[
 4\pi\chi(\Sigma)
 =\int_\Sigma 2K_k\,d\Area_k
 \ge c\Area_k(\Sigma)+\int_\Sigma|dU|_k^2\,d\Area_k>0.
\]
This formula is valid also for nonorientable surfaces, either directly
with densities or by passing to the orientation double cover.  The
classification of closed connected surfaces now gives precisely
$S^2$ and $\R P^2$.
\end{proof}

\begin{lemma}[Local exit estimate]\label{lem:stable-exit}
Let $\Sigma$ be a smooth two-sided stable surface satisfying
$H-\nu T=\mu$ in a Riemannian three-manifold $(X,h)$.  Suppose that
\eqref{eq:pressure-curvature-condition} holds on $\Sigma$.
Let $\Omega_0,\Omega_1$ be connected smooth domains with compact closures
in $\operatorname{int}\Sigma$ such that $\overline{\Omega_0}\subset\Omega_1$ and
$\partial\Omega_0\ne\varnothing$.  For every $x\in\Omega_0$ there
is a path from $x$ to $\partial\Omega_0$, contained in
$\overline{\Omega_0}$, whose length in $k=h|_\Sigma$ is at most
$2\pi/\sqrt c$.
\end{lemma}

\begin{proof}
Apply Lemma~\ref{lem:weighted-stability} in $\Omega_1$.  The resulting
$U$ is smooth on a neighborhood of $\overline{\Omega_0}$ and
satisfies $D(k,U)\ge c$ there.  In the compact Riemannian manifold
with boundary $\overline{\Omega_0}$, choose a path of least
$e^{-2U}k$-length from $x$ to $\partial\Omega_0$, and stop it at its
first boundary point.  Its interior is a geodesic in $\Omega_0$ and
its index form is nonnegative for compactly supported interior
variations.  Lemma~\ref{lem:conformal-length} gives the stated
$k$-length estimate.  This argument uses compactness only of the two
chosen domains, not completeness of $\Sigma$.
\end{proof}

For the original data $(M,g,t)$, weighted minimality means $\mu=0$.
Since $D(g,t)\ge1$, Proposition~\ref{prop:stable-surface} bounds every
closed connected two-sided stable weighted minimal surface by $2\pi$
and identifies its topology.  The same constant in
Lemma~\ref{lem:stable-exit} prevents an open stable sheet from extending
arbitrarily far inside a region where it has compact intersections with
ambient balls.
